\documentclass[10pt,a4paper]{amsart}
\usepackage{amsmath,amssymb,amsthm,graphicx}
% Retained preamble for the Stacks Project chapter snapshot in this directory.
% The chapter's own preamble.tex is a copy of the shared Stacks Project preamble; the block
% below is sliced MECHANICALLY, line for line, from it. Nothing was added, removed or
% redefined; the wrapper class, the xr-hyper/hyperref wiring and the externaldocument list
% are omitted because the assembler supplies the document class.
%
% lines taken from preamble.tex: 3, 5, 9, 11, 12, 156, 157, 158, ... (51 lines total)
%
% package coverage restored from the chapter's own preamble (2026-10-13): xy, xy-2cell

\usepackage[all]{xy}
\xyoption{2cell}
\UseAllTwocells

\newenvironment{reference}{\comment}{\endcomment}
\newenvironment{slogan}{\comment}{\endcomment}
\newenvironment{history}{\comment}{\endcomment}
\theoremstyle{plain}
\newtheorem{theorem}[subsection]{Theorem}
\newtheorem{proposition}[subsection]{Proposition}
\newtheorem{lemma}[subsection]{Lemma}
\theoremstyle{definition}
\newtheorem{definition}[subsection]{Definition}
\newtheorem{example}[subsection]{Example}
\newtheorem{exercise}[subsection]{Exercise}
\newtheorem{situation}[subsection]{Situation}
\theoremstyle{remark}
\newtheorem{remark}[subsection]{Remark}
\newtheorem{remarks}[subsection]{Remarks}
\numberwithin{equation}{subsection}
\def\lim{\mathop{\mathrm{lim}}\nolimits}
\def\colim{\mathop{\mathrm{colim}}\nolimits}
\def\Spec{\mathop{\mathrm{Spec}}}
\def\Hom{\mathop{\mathrm{Hom}}\nolimits}
\def\Ext{\mathop{\mathrm{Ext}}\nolimits}
\def\SheafHom{\mathop{\mathcal{H}\!\mathit{om}}\nolimits}
\def\SheafExt{\mathop{\mathcal{E}\!\mathit{xt}}\nolimits}
\def\Sch{\mathit{Sch}}
\def\Mor{\mathop{\mathrm{Mor}}\nolimits}
\def\Ob{\mathop{\mathrm{Ob}}\nolimits}
\def\Sh{\mathop{\mathit{Sh}}\nolimits}
\def\NL{\mathop{N\!L}\nolimits}
\def\CH{\mathop{\mathrm{CH}}\nolimits}
\def\proetale{{pro\text{-}\acute{e}tale}}
\def\etale{{\acute{e}tale}}
\def\QCoh{\mathit{QCoh}}
\def\Ker{\mathop{\mathrm{Ker}}}
\def\Im{\mathop{\mathrm{Im}}}
\def\Coker{\mathop{\mathrm{Coker}}}
\def\Coim{\mathop{\mathrm{Coim}}}
\DeclareMathSymbol{\boxtimes}{\mathbin}{AMSa}{"02}
\def\QCohstack{\mathcal{QC}\!\mathit{oh}}
\def\Cohstack{\mathcal{C}\!\mathit{oh}}
\def\Spacesstack{\mathcal{S}\!\mathit{paces}}
\def\Quotfunctor{\mathrm{Quot}}
\def\Hilbfunctor{\mathrm{Hilb}}
\def\Curvesstack{\mathcal{C}\!\mathit{urves}}
\def\Polarizedstack{\mathcal{P}\!\mathit{olarized}}
\def\Complexesstack{\mathcal{C}\!\mathit{omplexes}}
\def\Pic{\mathop{\mathrm{Pic}}\nolimits}
\def\Picardstack{\mathcal{P}\!\mathit{ic}}
\def\Picardfunctor{\mathrm{Pic}}
\def\Deformationcategory{\mathcal{D}\!\mathit{ef}}
\usepackage{cleveref}
\crefname{originallemma}{original lemma}{original lemmas}
\Crefname{originallemma}{Original lemma}{Original lemmas}
\makeatletter
\newlabel{categories-lemma-colimits-coproducts-coequalizers}{{002P}{source}{}{sourcepointer-1}{}}
\newlabel{categories-lemma-colimits-coproducts-coequalizers@cref}{{[originallemma][0][]002P}{[1][0][]source}{}{}{}}
\makeatother
\crefname{originallemma}{original lemma}{original lemmas}
\Crefname{originallemma}{Original lemma}{Original lemmas}
\makeatletter
\newlabel{categories-lemma-limits-products-equalizers}{{002N}{source}{}{sourcepointer-2}{}}
\newlabel{categories-lemma-limits-products-equalizers@cref}{{[originallemma][0][]002N}{[1][0][]source}{}{}{}}
\makeatother
\crefname{originallemma}{original lemma}{original lemmas}
\Crefname{originallemma}{Original lemma}{Original lemmas}
\makeatletter
\newlabel{definition-dgm}{{0FRI}{source}{}{sourcepointer-3}{}}
\newlabel{definition-dgm@cref}{{[originallemma][0][]0FRI}{[1][0][]source}{}{}{}}
\makeatother
\crefname{originallemma}{original lemma}{original lemmas}
\Crefname{originallemma}{Original lemma}{Original lemmas}
\makeatletter
\newlabel{lemma-gm-abelian}{{0FR0}{source}{}{sourcepointer-4}{}}
\newlabel{lemma-gm-abelian@cref}{{[originallemma][0][]0FR0}{[1][0][]source}{}{}{}}
\makeatother
\title{Stacks, Differential Graded Sheaves: lemma-dgm-abelian}
\author{The Stacks Project Authors}
\date{stacks-project:sdga.tex:local-snapshot-20261012; GNU FDL 1.2 or later}
\begin{document}
\maketitle
\noindent\textit{Editorial note.} Statement span is lines 1443-1464 of sources/stacks-sdga/pinned-original-source.tex; proof span is lines 1466-1509, the complete author proof (42 body lines).\par\medskip
\noindent\textit{Editorial note.} Navigation stubs: 4 cited-but-not-included labels are supplied as non-load-bearing pointers with their pinned tags where the snapshot tag index has them.\par\medskip
\noindent\textit{Editorial note.} extra\_packages is empty on purpose: the retained preamble already loads xy with options and a second declaration clashes in the assembler wrapper.\par\medskip
\noindent\textit{Editorial note.} Text is retained verbatim, including the author's own typos and placeholders.\par\medskip
\noindent\textit{Editorial note.} Queue index 14 of queue\_sdga.json (pb=1945, stmt 1443-1464, proof 1466-1509).\par\medskip
\noindent\textit{Editorial note.} Cross-chapter same-label warning (NOT an occupancy conflict, Chief ruling 2026-10-13 evening): label 'lemma-dgm-abelian' also exists in stacks-dga — a different proposition in a different chapter file.\par\medskip
\noindent\textit{Allowed standard prerequisites (editorial).}\begin{itemize}
\item Explicit prerequisite boundary. The cited inputs are prerequisites; the argument in the retained span is not replaced by any cited result.
\end{itemize}\medskip
\ifdefined\hypertarget\hypertarget{sourcepointer-1}{}\fi
\noindent\textit{Source reference.} Original lemma 002P: sources/stacks-sdga/pinned-original-source.tex, source lines 1509--1509. Cited by the retained statement or proof but not part of the collected spans; supplied only as a navigation pointer. The argument does not depend on its content.\par\medskip
\ifdefined\hypertarget\hypertarget{sourcepointer-2}{}\fi
\noindent\textit{Source reference.} Original lemma 002N: sources/stacks-sdga/pinned-original-source.tex, source lines 1509--1509. Cited by the retained statement or proof but not part of the collected spans; supplied only as a navigation pointer. The argument does not depend on its content.\par\medskip
\ifdefined\hypertarget\hypertarget{sourcepointer-3}{}\fi
\noindent\textit{Source reference.} Original lemma 0FRI: sources/stacks-sdga/pinned-original-source.tex, source lines 1396--1418. Cited by the retained statement or proof but not part of the collected spans; supplied only as a navigation pointer. The argument does not depend on its content.\par\medskip
\ifdefined\hypertarget\hypertarget{sourcepointer-4}{}\fi
\noindent\textit{Source reference.} Original lemma 0FR0: sources/stacks-sdga/pinned-original-source.tex, source lines 386--397. Cited by the retained statement or proof but not part of the collected spans; supplied only as a navigation pointer. The argument does not depend on its content.\par\medskip
\begin{lemma}
\label{lemma-dgm-abelian}
Let $(\mathcal{C}, \mathcal{O})$ be a ringed site.
Let $(\mathcal{A}, \text{d})$ be a differential graded $\mathcal{O}$-algebra.
The category $\textit{Mod}(\mathcal{A}, \text{d})$ is an abelian category
with the following properties
\begin{enumerate}
\item $\textit{Mod}(\mathcal{A}, \text{d})$ has arbitrary direct sums,
\item $\textit{Mod}(\mathcal{A}, \text{d})$ has arbitrary colimits,
\item filtered colimit in $\textit{Mod}(\mathcal{A}, \text{d})$ are exact,
\item $\textit{Mod}(\mathcal{A}, \text{d})$ has arbitrary products,
\item $\textit{Mod}(\mathcal{A}, \text{d})$ has arbitrary limits.
\end{enumerate}
The forgetful functor
$$
\textit{Mod}(\mathcal{A}, \text{d})
\longrightarrow
\textit{Mod}(\mathcal{A})
$$
sending a differential graded $\mathcal{A}$-module to its underlying
graded module commutes with all limits and colimits.
\end{lemma}

\begin{proof}
Let us denote
$F : \textit{Mod}(\mathcal{A}, \text{d}) \to \textit{Mod}(\mathcal{A})$
the functor in the statement of the lemma. Observe that
the category $\textit{Mod}(\mathcal{A})$ has properties (1) -- (5), see
Lemma \ref{lemma-gm-abelian}.

\medskip\noindent
Consider a homomorphism $f : \mathcal{M} \to \mathcal{N}$
of graded $\mathcal{A}$-modules. The kernel
and cokernel of $f$ as maps of graded $\mathcal{A}$-modules
are additionally endowed with differentials as in
Definition \ref{definition-dgm}. Hence these are also
the kernel and cokernel in $\textit{Mod}(\mathcal{A}, \text{d})$.
Thus $\textit{Mod}(\mathcal{A}, \text{d})$ is an abelian category
and taking kernels and cokernels commutes with $F$.

\medskip\noindent
To prove the existence of limits and colimits it is sufficient
to prove the existence of products and direct sums, see
Categories, Lemmas \ref{categories-lemma-limits-products-equalizers} and
\ref{categories-lemma-colimits-coproducts-coequalizers}.
The same lemmas show that
proving the commutation of limits and colimits with $F$
follows if $F$ commutes with direct sums and products.

\medskip\noindent
Let $\mathcal{M}_t$, $t \in T$ be a set of differential
graded $\mathcal{A}$-modules. Then we can consider the direct
sum $\bigoplus \mathcal{M}_t$ as a graded $\mathcal{A}$-module.
Since the direct sum of graded modules is done termwise, it is
clear that $\bigoplus \mathcal{M}_t$ comes endowed with a differential.
The reader easily verifies that this is a direct sum in
$\textit{Mod}(\mathcal{A}, \text{d})$. Similarly for products.

\medskip\noindent
Observe that $F$ is an exact functor and that
a complex $\mathcal{M}_1 \to \mathcal{M}_2 \to \mathcal{M}_3$
of $\textit{Mod}(\mathcal{A}, \text{d})$ is exact if and only if
$F(\mathcal{M}_1) \to F(\mathcal{M}_2) \to F(\mathcal{M}_3)$
is exact in $\textit{Mod}(\mathcal{A})$. Hence we conclude that (3)
holds as filtered colimits are exact in
$\textit{Mod}(\mathcal{A})$.
\end{proof}

\end{document}
